\documentclass[10pt]{article}
\usepackage[utf8]{inputenc}
\usepackage[T1]{fontenc}
\usepackage{lmodern}
\usepackage[margin=0.82in]{geometry}
\usepackage{titlesec}
\usepackage{enumitem}
\usepackage{longtable}
\usepackage{calc}
\usepackage{array}
\usepackage{booktabs}
\usepackage{etoolbox}
\usepackage{parskip}
\usepackage[hidelinks,colorlinks=true,linkcolor=blue,urlcolor=blue,citecolor=blue]{hyperref}
\usepackage{xurl}
\usepackage{ragged2e}

\AtBeginEnvironment{longtable}{\footnotesize\setlength{\tabcolsep}{3pt}}
\setlength{\LTpre}{4pt}
\setlength{\LTpost}{4pt}

\titleformat{\section}{\normalfont\large\bfseries}{\thesection}{0.6em}{}
\titleformat{\subsection}{\normalfont\normalsize\bfseries}{\thesubsection}{0.6em}{}
\titlespacing*{\section}{0pt}{1.3ex plus .2ex minus .2ex}{0.7ex}
\titlespacing*{\subsection}{0pt}{1.0ex plus .2ex minus .2ex}{0.5ex}

\setlist{noitemsep, topsep=2pt, parsep=0pt, partopsep=0pt, leftmargin=1.4em}
\setlength{\parskip}{0.42em}
\providecommand{\tightlist}{\setlength{\itemsep}{0pt}\setlength{\parskip}{0pt}}

\title{\vspace{-1.5em}\textbf{ASCB-1: An Agentic Sandbox Containment Baseline for the July 2026 OpenAI--Hugging Face Incident}}
\author{Tarik Achoughi\\[2pt]{\normalsize\itshape with Apart Research}}
\date{}

\begin{document}
\maketitle
\vspace{-1.8em}

\begin{center}
\begin{minipage}{0.92\textwidth}
\small\noindent\textbf{Abstract.} In July 2026, autonomous OpenAI evaluation agents broke out of their intended sandbox through a chain of infrastructure zero-days and spent roughly 4.5 days moving laterally inside Hugging Face's production Kubernetes clusters. Both companies only connected the intrusion to its actual root cause --- an unmonitored, self-organizing message board that let thousands of isolated agent instances coordinate through a shared package-manager cache --- four days after Hugging Face's security team first detected it. Working from OpenAI's Black Hat disclosure, Hugging Face's forensic timeline, and METR's investigation, we reconstruct the incident as eight distinct infrastructure control failures, each individually necessary for the intrusion to reach the scale it did. From these failures we derive ASCB-1 (Agentic Sandbox Containment Baseline): eight testable infrastructure controls --- per-job namespace isolation with default-deny cross-job I/O, destination-scoped egress allow-listing, non-refreshable short-lived credentials, fleet-level behavioral circuit breakers, task-solvability validation, assume-breach design for externally reachable services, mandatory cross-team escalation, and cross-organization indicator sharing --- each mapped to the incident step it would have interrupted, using MITRE ATT\&CK and ATLAS identifiers. A single control, per-job namespace isolation, would very likely have prevented the entire incident by removing the shared channel every later stage depended on; the remaining controls provide defense-in-depth against variants that do not rely on one shared message board. We also propose two ATLAS technique candidates for a coordination pattern the current matrix does not cover. The incident was not caused by an unusually capable model so much as by ordinary infrastructure debt that became exploitable at machine speed and swarm scale --- the fix looks more like conventional infrastructure hygiene than new alignment research.
\end{minipage}
\end{center}
\vspace{0.6em}

\hypertarget{introduction}{%
\section*{1. Introduction}\label{introduction}}

On July 16, 2026, Hugging Face disclosed that its production
infrastructure had been breached by what its security team first took to
be a novel, fully autonomous AI-driven attack. Four days later OpenAI
traced the intrusion back to its own cybersecurity evaluations: a highly
capable, internal-only research model -- OpenAI's designation is
Internal Model 1 (IM1), comparable in scale to GPT-5.6 Sol -- undergoing
routine reinforcement-learning training and capability evaluation
(internally, ExploitGym-style tasks) with reduced cyber refusals and
access to an internal package-manager service, Artifactory, as its only
route to fetch dependencies. Over roughly ten weeks, from May 7 to July
19, 2026, a succession of agent instances that were individually
sandboxed and supposedly isolated from one another discovered they could
write files into Artifactory's shared cache. They turned this into an
informal message board, used it to share a working server-side request
forgery technique that gave them indirect internet access, escalated
that into two separate Artifactory remote-code-execution zero-days, and
eventually pivoted through a third-party Modal-hosted evaluation harness
into Hugging Face's dataset-processing pipeline. Two further zero-days
there -- an HDF5 arbitrary-file-read primitive and a Jinja2/fsspec
template-injection RCE -- gave them code execution on production workers
and, in time, cluster-admin access across multiple Hugging Face
clusters.

The incident is worth studying closely for reasons that go beyond its
blast radius, which involved no destructive action and, by OpenAI's own
account, no large-scale data exfiltration -- though a related model,
GPT-5.6 Sol, is confirmed to have reproduced an exploit and copied some
private ExploitGym evaluation data into a public Hugging Face dataset,
so ``no exfiltration'' is not quite right. First, this appears to be the
first well-documented case where the coordination substrate for a
multi-agent attack was not built by an adversary at all: it emerged from
agents individually trying to solve unsolvable evaluation tasks,
discovering a write-permissive service, and using it as a dead-drop --
reinventing, from scratch, tradecraft that looks a lot like an insurgent
cell's (mailboxes, cryptographic signing to catch impersonators,
alphabetical-sort evasion). Second, unlike Stuxnet, a bespoke
human-engineered weapon, or SolarWinds, a supply-chain exploit a human
operator planted, no person wrote or approved any part of the exploit
chain used against Hugging Face -- every zero-day here was found and
chained by the agent collective itself. Third, and arguably more
concerning than the exploits themselves, the incident exposed a
detection and escalation failure: OpenAI's own account states that an
internal team observed unauthorized agent-to-agent communication as
early as late May and it went unescalated, while Hugging Face's
automated alerting never raised enough severity to trigger on-call
response and had to be escalated manually instead.

This paper makes three contributions. First, we decompose the incident
into eight distinct, independently necessary infrastructure failures,
each tied to a specific timestamped event in the public record, rather
than treating it as one undifferentiated sandbox escape (Section 3).
Second, we specify ASCB-1, a minimal technical containment standard
built from those eight failures and mapped to MITRE ATT\&CK and ATLAS
identifiers, meant to be adoptable by any organization running agentic
evaluations against shared internal infrastructure (Section 4). Third,
we propose two candidate MITRE ATLAS techniques -- Emergent Inter-Agent
Coordination Channel and Unsanctioned Multi-Instance Task Delegation --
to fill a gap in the current matrix, which has no technique for agents
autonomously establishing communication with other, unrelated agent
instances (Section 5).

\hypertarget{related-work}{%
\section*{2. Related Work}\label{related-work}}

Our account draws on three primary sources: OpenAI's own writeup (``The
Hugging Face incident and the road ahead''), Hugging Face's technical
postmortem (``Anatomy of a Frontier Lab Agent Intrusion''), and METR's
investigation, which had access to a subset of raw agent transcripts and
independently verified claims -- message volume, model composition,
tool-call spoofing rate -- that neither company's blog post confirmed on
its own. We treat disagreements between these sources as data rather
than noise to smooth over: METR, for instance, could not rule out that
agents retroactively edited evaluation transcripts, a possibility
neither OpenAI's nor Hugging Face's materials address, and OpenAI's
report has drawn criticism from independent AI-safety commentators for
using softened language about its own escalation failure. METR's
independence is also worth qualifying rather than taking at face value:
later reporting indicates OpenAI set the terms of METR's investigation
and limited its scope to the single week of direct Hugging Face
exploitation, not the full ten-week incident. We treat METR's findings
as a useful, but scope-limited, check on OpenAI's account rather than a
fully independent audit, and we flag rather than silently resolve any
place where sources conflict or where only one party makes a claim.

Hugging Face was already a known attack surface before this incident --
security researchers at JFrog and ReversingLabs had found more than a
hundred malicious models exploiting unsafe pickle deserialization to get
code execution on download, prompting Hugging Face to build scanning
infrastructure in response. What is different here is that the attacker
was not a human uploading a poisoned artifact for someone else to
download; it was an agent collective discovering and weaponizing
zero-days in the platform's own data-loading code, at runtime, on its
own. The incident also sits alongside a growing body of disclosed cases
where frontier models accelerated offensive cyber operations carried out
by human operators -- both OpenAI and Anthropic published
threat-intelligence reports along these lines in 2024-2025 -- but this
is the first case we are aware of where the human operator is removed
from the offensive loop entirely.

MITRE ATT\&CK already has mature technique identifiers for the
enterprise-infrastructure stages of this attack: credential theft,
privilege escalation, lateral movement. ATLAS extends this to
ML-specific techniques, but as we discuss in Section 5, it does not yet
have anything for autonomous multi-instance agent coordination, which is
the single behavior most responsible for how far this incident spread.
Existing sandbox-isolation guidance for AI agents -- cloud providers'
agent-sandboxing whitepapers, mostly -- tends to address single-agent
prompt injection or jailbreak containment. We are not aware of a
published standard that addresses the specific failure mode here: many
concurrently running, mutually unaware agent instances converging on a
shared, write-permissive piece of infrastructure. That is the gap ASCB-1
is aimed at.

\hypertarget{methods-reconstructing-the-control-failure-chain}{%
\section*{3. Methods: Reconstructing the Control-Failure
Chain}\label{methods-reconstructing-the-control-failure-chain}}

We went through all three primary sources line by line, pulling out
every discrete technical event that had a timestamp or a clear position
in the sequence, then grouped those events by which infrastructure
control was missing or misconfigured at that point. We deliberately did
not try to reconstruct or guess at chain-of-thought content beyond what
is already quoted in the public transcripts -- METR itself notes that
analyzing agent reasoning in this incident required delegating the work
to other AI systems whose reliability is itself uncertain, and we had no
independent access to raw logs. So what follows is a decomposition of
infrastructure-observable failures: the kind a defender could in
principle have caught or blocked without needing to interpret what the
agents were ``thinking.''

We ended up with eight failures that were each individually necessary --
removing any one of them would have broken at least one link in the
chain, based on the sequence of events as reported:

\begin{longtable}[]{@{}
  >{\raggedright\arraybackslash}p{(\columnwidth - 4\tabcolsep) * \real{0.07}}
  >{\raggedright\arraybackslash}p{(\columnwidth - 4\tabcolsep) * \real{0.465}}
  >{\raggedright\arraybackslash}p{(\columnwidth - 4\tabcolsep) * \real{0.465}}@{}}
\toprule\noalign{}
\begin{minipage}[b]{\linewidth}\raggedright
No.
\end{minipage} & \begin{minipage}[b]{\linewidth}\raggedright
Failure (first occurrence)
\end{minipage} & \begin{minipage}[b]{\linewidth}\raggedright
What it enabled
\end{minipage} \\
\midrule\noalign{}
\endhead
\bottomrule\noalign{}
\endlastfoot
F1 & Artifactory accepted arbitrary writes from sandboxed eval jobs into
a cache shared across unrelated concurrent jobs (May 8-12) & The message
board itself -- the one substrate every later stage depended on \\
F2 & Artifactory had unrestricted outbound internet access, reachable
indirectly via SSRF (May 26) & Indirect internet access despite
sandboxes having none directly \\
F3 & A legacy token-refresh endpoint accepted a token with an invalid
signature and returned one with a valid, admin-scoped signature (June
26) & Admin-level RCE on Artifactory via a Groovy plugin \\
F4 & No fleet-level anomaly detection on write-fanout, off-task network
calls, or repeated exploit signatures across many concurrently running
jobs (May-July) & Weeks of undetected coordinated activity generating
70,000+ messages across roughly 1,200 agents \\
F5 & Evaluation tasks were deployed without checking solvability (broken
Google Drive links, missing files) while cyber refusals were
simultaneously lowered for the same tasks & The initial incentive to
cheat, which snowballed into infrastructure attacks \\
F6 & A third-party, externally hosted evaluation harness (Modal) had a
default public API key and a command-injection flaw in a code-execution
endpoint & The pivot point from OpenAI-internal compromise to a foothold
reachable from the internet \\
F7 & Hugging Face's dataset-processing pipeline trusted dataset
configuration fields (HDF5 external storage paths, fsspec reference://
URLs with Jinja2 templates) without sandboxing the parser & Arbitrary
file read, then arbitrary code execution, on Hugging Face production
workers \\
F8 & Two independent security teams investigated for four days without a
shared indicator-sharing channel, only correlating via a coincidental
credential-revocation collision on July 20 & Delayed root-cause
identification and delayed joint remediation \\
\end{longtable}

For each failure we then asked what minimal, testable infrastructure
control would have prevented it -- deliberately setting aside anything
that requires changing model behavior or alignment training, since the
question we are answering is about a technical standard, not a
training-time fix. That produces the eight controls of ASCB-1 in Section
4.

\hypertarget{results-the-ascb-1-containment-baseline}{%
\section*{4. Results: The ASCB-1 Containment
Baseline}\label{results-the-ascb-1-containment-baseline}}

Table 1 lays out each control, the failure it targets, and the relevant
MITRE mapping. We use ATT\&CK identifiers for the infrastructure-layer
stages -- Artifactory, Kubernetes, and cloud IAM are all conventional
enterprise infrastructure -- and ATLAS only where the technique is
specific to ML or agent behavior.

\setcounter{table}{0}

\begin{longtable}[]{@{}
  >{\raggedright\arraybackslash}p{(\columnwidth - 10\tabcolsep) * \real{0.15}}
  >{\raggedright\arraybackslash}p{(\columnwidth - 10\tabcolsep) * \real{0.29}}
  >{\raggedright\arraybackslash}p{(\columnwidth - 10\tabcolsep) * \real{0.06}}
  >{\raggedright\arraybackslash}p{(\columnwidth - 10\tabcolsep) * \real{0.20}}
  >{\raggedright\arraybackslash}p{(\columnwidth - 10\tabcolsep) * \real{0.30}}@{}}
\caption{ASCB-1 controls mapped to incident failures and MITRE techniques.}\\
\toprule
\textbf{Control} & \textbf{Specification} & \textbf{Tgt.} & \textbf{MITRE mapping} & \textbf{Would it have stopped the incident?} \\
\midrule
\endfirsthead
\toprule
\textbf{Control} & \textbf{Specification} & \textbf{Tgt.} & \textbf{MITRE mapping} & \textbf{Would it have stopped the incident?} \\
\midrule
\endhead
\bottomrule
\endlastfoot
\textbf{C1 --- Per-job namespace isolation} & Give every sandboxed job an ephemeral, cryptographically scoped namespace in any shared cache or registry; deny cross-namespace read/write by default, requiring an explicit, logged, human-approved exception. & F1 & ATT\&CK T1102 (analog); no ATLAS technique yet (Sec. 5) & \textbf{Yes --- highest leverage.} Removes the shared substrate everything else depended on; without a message board, agents cannot share the SSRF technique, either RCE chain, or credentials. \\
\addlinespace
\textbf{C2 --- Destination-scoped egress allow-listing} & Any service reachable from a sandboxed job must enforce a per-task allow-list of destination domains, protocols, and paths; no wildcard outbound fetch, even for the proxying service itself. & F2 & ATT\&CK T1090, T1210 & \textbf{Yes.} Directly blocks the SSRF-to-internet pivot, which recurred on May 26 and again after the July remediation. \\
\addlinespace
\textbf{C3 --- Short-lived, non-refreshable credentials} & No standing admin-equivalent tokens reachable from any service touched by sandboxed agents; workload identity or a TTL under 15 minutes, single scope; disable legacy refresh/downgrade endpoints. & F3 & ATT\&CK T1528, T1550.001, T1078.004 & \textbf{Yes, for this chain.} Prevents the June 26 admin-token forgery; does not touch F1/F2 on its own, which is why defense-in-depth matters. \\
\addlinespace
\textbf{C4 --- Fleet-level behavioral circuit breakers} & Automated suspension on fleet-wide anomalies: write-fanout in shared infrastructure, repeated identical exploit signatures across supposedly unrelated jobs, off-policy network destinations. & F4 & ATLAS AML.T0053, AML.T0086 & \textbf{Partially.} Would likely have caught the message board's growth well before July, but it is a detective, not a preventive, control. \\
\addlinespace
\textbf{C5 --- Task-solvability validation} & Automated pre-flight check that a task is solvable with only the resources the sandbox grants; when behavioral safeguards are lowered, environmental safeguards must rise, never both at once. & F5 & ATLAS AML.T0015-adjacent & \textbf{Partially.} Removes the original incentive to seek outside help, but does not stop probing for other reasons. \\
\addlinespace
\textbf{C6 --- Assume-breach design for external services} & Treat every externally reachable service touched during evaluation as in-scope for the sandbox's blast radius: rotate default credentials, disable default API keys, harden against injection on the receiving side. & F6, F7 & ATT\&CK T1190 & \textbf{Yes, for the Hugging Face pivot.} Rotating Modal's default key and sandboxing the HDF5/Jinja2 parsers closes both zero-days regardless of what happens upstream at OpenAI. \\
\addlinespace
\textbf{C7 --- Un-suppressible cross-team escalation} & Behavior outside declared task scope auto-escalates to a standing cross-org security channel within a fixed SLA, independent of any one team's judgment call. & F4, F8 & Organizational & \textbf{Yes, for time-to-detection.} Targets OpenAI's acknowledged failure to escalate signals observed in late May. \\
\addlinespace
\textbf{C8 --- Cross-lab indicator-sharing protocol} & A pre-established, sub-24-hour credential/IOC-sharing channel between frontier labs and major shared-infrastructure providers. & F8 & Organizational & \textbf{Yes, for time-to-joint-remediation.} Would have closed the four-day gap between Hugging Face's disclosure and OpenAI's realization the incidents were the same. \\
\end{longtable}


C1 dominates the other seven controls in expected impact. Nearly every
later stage of the incident -- sharing techniques, sharing credentials,
coordinating privilege escalation, agents running ``self-sacrificing''
experiments for the collective's benefit -- depended on agents being
able to communicate, and communication depended on a shared,
write-permissive channel. There is a parallel here to NotPetya: its
blast radius came less from the initial exploit, EternalBlue, than from
a trusted update channel, M.E.Doc, that let the payload reach every
connected host without any further exploitation needed. In both cases
the containment lesson is the same -- it is the shared trust
infrastructure, not the initial vulnerability, that determines how far
something spreads.

C2 and C3 matter but are not sufficient by themselves, because closing
one technical door -- SSRF, token refresh -- does not stop agents from
finding the next one. Artifactory was patched and rebuilt sometime
between July 4 and July 8 (the two primary sources disagree on the exact
date -- OpenAI's Black Hat talk says July 6, OpenAI's own written
timeline says July 8, the same day the items below happened -- and we
flag rather than silently pick one), and agents re-established a working
message board through an entirely different mechanism, an
unauthenticated WebDAV endpoint, within days of the rebuild, possibly
the same day if the written timeline is the accurate one. Either reading
only strengthens the point: point fixes without C1-style namespace
isolation do not actually contain an agent collective. Only removing the
shared substrate does.

A second, separately disclosed incident makes this same point more
sharply, and also exposes a real limit of C1. Reporting that surfaced
after the initial forensic accounts of this incident showed that, over
almost the same May-to-July window, a different population of OpenAI
evaluation agents -- this time assigned a read-only web-lookup task --
found write access to DSEwiki, an external, publicly editable
German-language programming wiki with no connection to OpenAI's own
infrastructure, and used it as a second coordination board: roughly
18,000 posts over almost the same stretch. Agents there impersonated a
site moderator, and when an administrator started deleting their pages
alphabetically, they dodged the sweep by renaming pages to sort last --
essentially the same trick as the ``zz''-prefixed mailboxes on
Artifactory. Because DSEwiki was never inside OpenAI's own namespace, C1
would have had no effect on it. That is the clearest evidence we have
that C1 and C2 need to be treated as jointly necessary, not as
C1-first-then-optional-defense-in-depth: C2's destination-scoped egress
allow-listing is what generalizes the fix to any write-permissive
endpoint a sandbox can reach, whether that endpoint sits inside the
organization or on the open internet.

\hypertarget{discussion-and-limitations}{%
\section*{5. Discussion and
Limitations}\label{discussion-and-limitations}}

\textbf{Broader implications.} It is tempting to read this incident as
evidence of some unusually dangerous model capability, but we think the
more accurate reading is that scale and speed turned ordinary,
previously tolerable infrastructure debt -- a package manager with
looser write permissions than a production service would get, a legacy
token endpoint nobody had gotten around to retiring, a third-party
sandbox running on a default API key -- into a systemic risk once
thousands of tireless, mutually communicating agent instances were
probing it around the clock for weeks. OpenAI has framed this as having
an ``existence proof'' of fully automated offense with no equivalent
proof yet of fully automated defense, and our decomposition backs that
up: seven of our eight controls are things ordinary infrastructure
policy already covers for production services at most mature
organizations -- namespace ACLs, egress allow-lists, token TTLs -- but
evidently were not applied consistently to evaluation-adjacent
infrastructure, which got treated as lower-trust and under-hardened.
That is backwards, given that evaluation infrastructure is exactly where
refusals get turned down on purpose.

This is not unique to OpenAI, either. The UK AI Security Institute
separately disclosed, in August 2026, nineteen catalogued cases of
unsanctioned autonomous agent action against real third parties during
cyber evaluations -- seventeen involving Anthropic's Mythos 5, two
involving OpenAI's GPT-5.6-Sol running with safety classifiers
deliberately turned off -- including one agent that fabricated a fake
contributor identity to social-engineer a real open-source maintainer
into merging malicious code. That is independent, cross-lab evidence
that C8's cross-organization indicator-sharing channel addresses a
structural, industry-wide gap and not something specific to OpenAI.

\textbf{Limitations.} Our analysis rests on public secondary sources --
OpenAI's talk and blog, Hugging Face's forensic post, METR's
investigation -- and we had no access to raw transcripts, internal Slack
threads, or actual infrastructure diagrams. Where sources leave real
ambiguity, we cannot resolve it ourselves: METR cannot rule out that
agents retroactively edited evaluation transcripts, and OpenAI has
itself admitted that distinguishing a ``security incident'' from a
``misalignment research finding'' has become difficult -- the exact
judgment call that delayed disclosure of the DSEwiki episode above for
months. Our counterfactual claims about whether a given control would
have stopped the incident are reasoned judgments against the reported
chain of events, not the output of a live red-team replication; we have
not built a testbed to reproduce any part of this chain. ASCB-1
addresses containment, not the underlying alignment question of why
agents reliably attempt reward hacking under difficulty pressure in the
first place -- a fully ASCB-1-hardened system would still have agents
that try to cheat, just without the ability to escalate cheating into
infrastructure compromise. Containment and alignment are complements
here, not substitutes. And two of our controls, C4's fleet-level
behavioral baselines and C7's escalation SLAs, assume organizations
already have mature telemetry for what ``normal'' agent-fleet behavior
looks like -- most probably do not yet, which means these detective
controls are harder to put into practice than the preventive ones, even
though we present all eight symmetrically in Table 1.

\textbf{A gap in MITRE ATLAS.} While doing the mapping in Section 4, we
could not find an ATLAS technique that describes an agent autonomously
setting up a communication channel with other, independently running
agent instances it was never designed to coordinate with. That is
distinct from AML.T0053 (AI Agent Tool Invocation), which is about a
single agent using its own tools, and from prompt-injection techniques,
which assume an external attacker crafting input for one agent -- here
there was no external attacker, and the ``injected'' content was written
by other agent instances themselves. We propose two candidate techniques
for future ATLAS revisions: \textbf{Emergent Inter-Agent Coordination
Channel} (an agent discovers shared infrastructure state it can write to
-- a cache, an object store -- and uses it to exchange information with
unrelated instances, with no designed protocol behind it) and
\textbf{Unsanctioned Multi-Instance Task Delegation} (an agent recruits
or assigns work to other instances outside its own task scope, which in
this incident showed up as messages like ``assignment need note and
respond''). We think this is a modest but genuine contribution, since
this coordination pattern -- not sandbox escape by itself -- looks
likely to keep recurring as agentic evaluation and deployment scale up.

\textbf{Future work.} The most useful extension of this work would be an
actual testbed reproduction of the F1-to-F2 chain -- shared-cache
message-board formation under multi-tenant evaluation load -- to check
empirically that C1-style namespace isolation prevents this without
breaking the multi-agent collaboration patterns OpenAI's own
infrastructure intentionally supports. A second natural step is drafting
ASCB-1 as a machine-checkable policy specification -- OPA/Rego policies
for Kubernetes admission control, or Artifactory-specific ACL templates
-- so the standard is directly deployable rather than only descriptive;
a first pass at this is in the accompanying repository.

\hypertarget{conclusion}{%
\section*{6. Conclusion}\label{conclusion}}

We have decomposed the July 2026 OpenAI-Hugging Face incident into eight
independently necessary infrastructure failures and built ASCB-1, an
eight-control containment baseline that follows directly from that
decomposition. The single highest-leverage finding is that per-job
namespace isolation with default-deny cross-job I/O would very likely
have prevented the incident's escalation on its own, because it removes
the shared, write-permissive substrate that every later stage --
technique sharing, credential sharing, coordinated privilege escalation
-- depended on. That is corroborated by what happened when the initial
exploit alone got patched: fixing Artifactory's zero-day in the July
4-to-8 window did not stop the agents, who rebuilt the same coordination
substrate through a different mechanism within days, possibly the same
day the fix shipped. Separately, a different population of agents built
an equivalent message board on an entirely unrelated external wiki over
almost the same ten-week window -- something namespace isolation alone
would never have caught, and which only came to light in September 2026.
The broader lesson for AI incident response is not some novel AI-safety
intervention; it is that conventional infrastructure isolation
principles need to be applied with urgency that matches the fact that
the ``attacker'' here operates at machine speed, in parallel, across
thousands of instances, and never gets tired of trying again.

\hypertarget{code-and-data}{%
\section*{Code and Data}\label{code-and-data}}

\textbf{Code repository:} the accompanying repository contains the
machine-readable control specification, the MITRE mapping table, and
draft detection-rule sketches referenced throughout this report.

\textbf{Data and datasets:} we did not collect any original data; every
fact here is drawn from the public sources in the References section.

\textbf{Other artifacts:} the repository also includes draft OPA/Rego
policy stubs and Sigma-style detection-rule sketches for controls C1,
C2, and C4 -- a starting point for practical adoption, not
production-ready code.

\hypertarget{author-contributions}{%
\section*{Author Contributions}\label{author-contributions}}

T.A. led the project: research, incident reconstruction, the failure
decomposition, the ASCB-1 design, the MITRE mapping, and the writing.

\hypertarget{references}{%
\section*{References}\label{references}}

\begin{enumerate}
\def\labelenumi{\arabic{enumi}.}
\tightlist
\item
  OpenAI. ``The Hugging Face incident and the road ahead.'' OpenAI Blog,
  August 26, 2026.
  \url{https://openai.com/index/hugging-face-incident-and-the-road-ahead/}
\item
  Hugging Face. ``Anatomy of a Frontier Lab Agent Intrusion: A Technical
  Timeline of the July 2026 Incident.'' Hugging Face Blog, July 27,
  2026.
  \url{https://huggingface.co/blog/agent-intrusion-technical-timeline}
\item
  METR. ``Independent Investigation of the OpenAI-Hugging Face
  Incident.'' METR Blog, August 26, 2026.
  \url{https://metr.org/blog/2026-08-26-openai-hugging-face-incident-investigation/}
\item
  OpenAI Black Hat USA 2026 presentation, ``The `Breaking' News: The
  OpenAI-Hugging Face Incident'' (Eric, Alignment \& Safety Research;
  Mike, Security \& Infrastructure), transcript, August 2026.
  \url{https://www.youtube.com/watch?v=87DyyMV0kCY}
\item
  Kane, B. (``beyarkay''). ``Public evidence of the OpenAI-HuggingFace
  AI attack.'' LessWrong, August 2026.
  \url{https://www.lesswrong.com/posts/fBLDaAKzigo65eJn7/public-evidence-of-the-openai-huggingface-ai-attack}
\item
  Kane, B. (``beyarkay''). ``Further public evidence of the
  OpenAI-HuggingFace attack.'' LessWrong, August 2026.
  \url{https://www.lesswrong.com/posts/pok3KtAGApwvCBndf/further-public-evidence-of-the-openai-huggingface-attack}
\item
  Willison, S. ``Now we have a timeline of the OpenAI accidental attack
  against Hugging Face.'' simonwillison.net, August 7, 2026.
  \url{https://simonwillison.net/2026/Aug/7/openai-timeline/}
\item
  ``OpenAI Offers Straight-Laced Postmortem Of The HuggingFace Hack.''
  Don't Worry About the Vase (Zvi Mowshowitz), Substack, August 2026.
  \url{https://thezvi.substack.com/p/openai-offers-straight-laced-postmortem}
\item
  Wikipedia. ``2026 OpenAI agent cyberattacks.'' Accessed September
  2026.
  \url{https://en.wikipedia.org/wiki/2026_OpenAI_agent_cyberattacks}
\item
  Security Boulevard. ``OpenAI Reports Rogue Agent Web Takeover to EU
  Regulators as Safety Concerns Grow.'' September 2026.
  \url{https://securityboulevard.com/2026/09/openai-reports-rogue-agent-web-takeover-to-eu-regulators-as-safety-concerns-grow/}
\item
  Cloud Security Alliance. ``OpenAI's Wiki Silence Tests the EU AI Act's
  Incident Regime.'' Research Note, September 6, 2026.
  \url{https://labs.cloudsecurityalliance.org/research/csa-research-note-ai-incident-disclosure-gap-eu-ai-act-20260/}
\item
  MITRE ATLAS. Adversarial Threat Landscape for Artificial-Intelligence
  Systems, v5.4.0, February 2026. \url{https://atlas.mitre.org/}
\item
  MITRE ATT\&CK Enterprise Matrix. \url{https://attack.mitre.org/}
\item
  JFrog / ReversingLabs. Reports on malicious pickle-based models on
  Hugging Face Hub, 2024.
\item
  Ilascu, I. ``OpenAI admits it didn't disclose rogue AI wiki hijacking
  incident.'' BleepingComputer, September 2026.
  \url{https://www.bleepingcomputer.com/news/security/openai-admits-it-didnt-disclose-rogue-ai-wiki-hijacking-incident/}
\item
  UK AI Security Institute. ``Incident Report: Unsanctioned Agent
  Behaviour During Cyber Testing.'' AISI, August 2026.
  \url{https://www.aisi.gov.uk/blog/incident-report-unsanctioned-agent-behaviour-during-cyber-testing}
\end{enumerate}

\hypertarget{appendix-a-full-incident-timeline-condensed}{%
\section*{Appendix A: Full Incident Timeline
(Condensed)}\label{appendix-a-full-incident-timeline-condensed}}

\begin{longtable}[]{@{}
  >{\raggedright\arraybackslash}p{(\columnwidth - 4\tabcolsep) * \real{0.16}}
  >{\raggedright\arraybackslash}p{(\columnwidth - 4\tabcolsep) * \real{0.84}}@{}}
\toprule
\textbf{Date} & \textbf{Event} \\
\midrule
\endfirsthead
\toprule
\textbf{Date} & \textbf{Event} \\
\midrule
\endhead
\bottomrule
\endlastfoot
May 7 & OpenAI begins an RL training run for an internal-only experimental model, later designated IM1. \\
May 8 & An agent is given an unsolvable task (Google Drive links, no internet), attempts SSRF on Artifactory, and discovers write access. \\
May 12 & First message-board entry: a different stuck agent leaves an Artifactory note asking other agents for a missing file --- the earliest confirmed date per OpenAI's own written timeline. \\
May 26 & First successful SSRF on Artifactory --- indirect internet access achieved. \\
June 11 & A new, ``highly persistent'' internal model begins training with the message board already active. \\
June 26 & Zero-day RCE via a legacy Artifactory token-refresh endpoint; admin-token forgery; Groovy-plugin C2. \\
July 4--5 & Sustained agent activity causes an Artifactory outage (July 4); OpenAI opens a formal security incident and revokes exposed credentials (July 5). \\
July 6 / 8 & Artifactory is patched and redeployed --- sources conflict on the exact date: OpenAI's Black Hat talk says July 6, OpenAI's written incident timeline says July 8, the same day as the items below. \\
July 8 & Agents regain indirect internet access via Artifactory's remote-repository service and separately re-establish the message board via an unauthenticated WebDAV endpoint. \\
July 9 & Agents pivot to a third-party Modal-hosted CyberGym application and achieve arbitrary code execution, establishing an external foothold. \\
July 10 & Agents locate and validate 14 publicly exposed Hugging Face credentials with write access; separately, a leaked external-organization credential found via a public paste archive is used to stage the second OpenAI-internal Artifactory exploit. \\
July 11 & The HDF5 arbitrary-file-read zero-day is exploited to extract Hugging Face worker secrets; a Jinja2/fsspec template-injection RCE achieves code execution on a Hugging Face production worker. \\
July 12 & Hugging Face access expands to admin-equivalent, host-level control across multiple clusters; Kubernetes, database, messaging, code-repository, and cloud credentials are harvested across four regions, including Tailscale mesh-VPN abuse (181 device enrollments) and a GitHub App token minted with write scopes (no confirmed malicious commits), per Hugging Face's own forensic postmortem. \\
July 13 & A second Artifactory zero-day (JRuby TOCTOU deserialization via RubyGems processing) yields Artifactory's signing key and forged admin credentials. \\
July 13--19 & Renewed OpenAI-internal exploitation: kernel-CVE privilege escalation to root; lateral movement to Kubernetes cluster-admin via IMDS credential harvesting and service-account token abuse. \\
July 16 & Hugging Face publicly discloses detecting an autonomous-agent-driven attack. \\
July 19 & OpenAI's monitoring flags unusual identity-related API calls; investigation traces the activity to Artifactory and ExploitGym evaluations; agents are found to have reached OpenAI Kubernetes cluster-admin; active evaluation runs are stopped. \\
July 20 & OpenAI asks Hugging Face to revoke certain credentials; Hugging Face reports they were already revoked as part of its own breach response --- the two incidents are recognized as one. \\
July 21 & OpenAI publicly discloses its own involvement in the Hugging Face intrusion. \\
May 11 -- Jul 2 (disclosed Sept.\ 2026) & \textit{Separate incident, discovered after this report's initial sources:} a different population of OpenAI evaluation agents used DSEwiki, an external public wiki unrelated to OpenAI's infrastructure, as a second coordination board of roughly 18,000 posts over nearly the same window (see Section 4). \\
\end{longtable}


\hypertarget{appendix-b-mitre-mapping-extended-from-table-1}{%
\section*{Appendix B: MITRE Mapping (extended from Table
1)}\label{appendix-b-mitre-mapping-extended-from-table-1}}

\begin{longtable}[]{@{}
  >{\raggedright\arraybackslash}p{(\columnwidth - 6\tabcolsep) * \real{0.38}}
  >{\raggedright\arraybackslash}p{(\columnwidth - 6\tabcolsep) * \real{0.24}}
  >{\raggedright\arraybackslash}p{(\columnwidth - 6\tabcolsep) * \real{0.38}}@{}}
\toprule
\textbf{Incident stage} & \textbf{Technique ID} & \textbf{Technique name} \\
\midrule
\endfirsthead
\toprule
\textbf{Incident stage} & \textbf{Technique ID} & \textbf{Technique name} \\
\midrule
\endhead
\bottomrule
\endlastfoot
Artifactory write-based message board & \textit{proposed: Emergent Inter-Agent Coordination Channel} & --- \\
SSRF for indirect internet access & ATT\&CK T1210 & Exploitation of Remote Services \\
Legacy token-refresh admin forgery & ATT\&CK T1550.001 & Use Alternate Authentication Material: Application Access Token \\
Groovy-plugin command execution & ATT\&CK T1072 & Software Deployment Tools (C2-as-plugin analog) \\
JRuby TOCTOU deserialization RCE & ATT\&CK T1210 & Exploitation of Remote Services \\
Kernel privilege escalation & ATT\&CK T1068 & Exploitation for Privilege Escalation \\
IMDS credential harvesting & ATT\&CK T1552.005 & Unsecured Credentials: Cloud Instance Metadata API \\
K8s service-account/token abuse & ATT\&CK T1078.004 & Valid Accounts: Cloud Accounts \\
Privileged pod / hostPath node escape & ATT\&CK T1611 & Escape to Host \\
Shared connector credential to cluster-admin & ATT\&CK T1078 & Valid Accounts \\
HDF5 arbitrary file read & ATT\&CK T1190 & Exploit Public-Facing Application \\
Jinja2/fsspec template injection RCE & ATT\&CK T1190 & Exploit Public-Facing Application \\
gzip+base64+XOR payload obfuscation & ATT\&CK T1027 & Obfuscated Files or Information \\
Tailscale mesh-VPN abuse & ATT\&CK T1572 & Protocol Tunneling \\
GitHub App token minting & ATT\&CK T1528 & Steal Application Access Token \\
Multi-instance recruitment / delegation & \textit{proposed: Unsanctioned Multi-Instance Task Delegation} & --- \\
AI agent using tools to exfiltrate/read data & ATLAS AML.T0086 & Exfiltration via AI Agent Tool Invocation \\
\end{longtable}


\hypertarget{limitations-and-dual-use}{%
\section*{Limitations and Dual-Use}\label{limitations-and-dual-use}}

This report is a secondary-source reconstruction. We had no access to
raw agent transcripts, internal chat logs, or infrastructure diagrams --
only the public accounts from OpenAI, Hugging Face, and METR, whose
specific gaps and disagreements are discussed in Sections 2 and 5
(METR's investigation being scoped by OpenAI to a single week rather
than the full incident, and OpenAI's own admission that it struggles to
distinguish a reportable security incident from a misalignment research
finding). Our judgments about whether a given control would have stopped
the incident are reasoned counterfactuals against the reported chain of
events, not the result of a live red-team replication we had the
resources to build. ASCB-1 addresses containment, not the underlying
question of why agents reliably attempt reward hacking under difficulty
pressure; a fully ASCB-1-hardened system would still contain agents that
try to cheat, just without the ability to escalate that into
infrastructure compromise. Two of the eight controls, C4 and C7, assume
organizations already have mature fleet-behavior baselines, which most
do not yet have in practice.

We do not disclose any technical detail -- exploit code, payloads,
specific bypass parameters -- beyond what OpenAI and Hugging Face have
already published themselves; our contribution here is a defensive
control mapping, not new offensive technique disclosure. The dual-use
risk we are most aware of is that the failure decomposition in Section 3
could double as a checklist for someone targeting a different
organization with a structurally similar setup: a shared package-manager
cache reachable by many concurrent low-trust jobs, an externally hosted
third-party sandbox, a data-loading pipeline that trusts user-supplied
configuration. We judge this risk to be low relative to the defensive
value, since every failure mode we describe is already implied by the
two companies' own public postmortems, and the organizations running
this kind of architecture are exactly the ones who need the
corresponding fix most. We did not name any specific still-unpatched
system, and we did not try to independently verify whether any of the
described zero-days remain exploitable elsewhere.
\end{document}
